\section{Metrics}
\subsection{Introduction}
As mentioned befor, there are a plethora of different methods to detect outliers.
To make the results comparable, we need to use performance metrics.
The slection of appropiate metrics is a difficult but crucial task, to make sure that the results are meaningful and comparable.
On the one hand, to make methods comparable in research, it is necessary to introduce metrics that are widely used and accepted by the research community.
On the other hand, the metrics should be suitable for the specific dataset and the specific problem of outlier detection, to make sure that the results are meaningful and not misleading.
Therefor I will begin with metrics, that are accepted by the research community, and end with a metric that is suitable for the specific datasets I use in this thesis. 
The metrics accepted by the research community are: the true positive rate and the false positive rate.
After that I will clarify more sophisticated metrics, such as the F1-Score, Accuracy and Precision.
In the end I will explain the h-measure and decribe why it is suitable for the specific datatests I will use for this thesis.
\todo{adding: concentrating on binary classification, not multiclass or hierarchical classification, because the datasets I use are binary.}
\todo{adding: explaining why wont use runtime}
\todo{critique other paper wich use not suitabl metrics, }
\todo{explain that a whole framework of metrics has to be implemented, because every metric has its own advantages and disadvantages.}
\textcolor{red}{Supervised Machine Learning (ML) has several ways of evaluating the performance of
learning algorithms and the classifiers they produce. Measures of the quality of clas-
sification are built from a confusion matrix which records correctly and incorrectly
recognized examples for each class. This paper argues that the measures commonly used now (accuracy, precision, recall,
F-Score and ROC Analysis) do not fully meet the needs of learning problems in which the classes are 
equally important and where several algorithms are compared. Our find-
ings agree with those of [1] who surveys the comparison of algorithms on multiple data
sets. His survey, based on the papers published at the International Conferences on ML
2003–2004, notes that algorithms are mainly compared on accuracy.
The vast majority of ML research focus on the settings where the examples are assumed
to be identically and independently distributed (IID)}

\subsection{implement to describe wich metric is commonly used}
\textcolor{green}
{We found 276 papers that clearly
presented the performance metrics of their proposed models.
Figure 6 shows that the performance metric used most was
True Positive Rate (TPR), which is also known as detection
date, sensitivity, and recall. It measures the anomalies that are
correctly classified. Moreover, 116 papers used False Positive
Rate (FPR) as a performance metric. This metric measures
anomalies that are falsely classified, and it can be known
as false alarm rate as well.  In addition, we find that 64 of the
290 papers used only one performance metric, and most of
those papers used only accuracy or AUC, which is not suffi-
cient to determine the quality performance of the ML model.
Furthermore, a lot of papers present computational
performance metrics in addition to performance metrics, such
as CPU utilization, execution time, training time, testing time,
Those two metrics are widely used in research. 
In the paper \citep[postfix]{nassif2021}, the authers found that the true positive rate is widley used, while the false positive rate is used in only 116 of the 290 papers.
and computational time \citep[postfix]{nassif2021}}


\subsection{Traditional Metrics - Confusion Matrix}
A simple way to evaluate the performance of a method, is to count the rightly and wrongly assigned instances to the classes.
In supervised outlier detection, the goal is to classify the presented data into outliers and inliers, using the ground truth of labled datasets.
Comparing it with the ground truth, you can plot them into a matrix, resulting in the confusion matrix.
The horizontal presents the true class and the vertical presents the hypothesized class of the method.

%\input{Images/Confusion Matrix.pdf}

The diagonal presenting the results the methods made, that are correct.
Outliers that are correctly classified by the method are called true positives, while inliers that are correctly classified are called true negatives.
The other diagonal presents the falsly classified data instances.
Outliers that are wrongly classified as inliers are called false negatives, while inliers that are wrongly classified as outliers are called false positives.
All basic metrics are derived fromt the confusion matrix \citep[p.4]{saito2015precision}. 
The most often used metrics derived from the confusion matrix are:

\subsection{Defition of Accuracy, Precision, True Positive Rate and Specificity}
\enumerate

    \item Accuracy: Measures the overall efectiveness of the method, by calculating the proportion of correctly classified instances to the total number of instances.
    It is defined as:
    
    \begin{equation}
        \text{Accuracy} = \frac{TP + TN}{TP + FN + FP + TN}
    \end{equation}

    \item Precision: Measures the proportion of true positives to the amount of all instances that are predicted as positives.
    The focus is the class agreement of the data labels with the positives labels given by the classifier.
    It is defined as: 

    \begin{equation}
        \text{Precision} = \frac{TP}{TP + FP}
    \end{equation}

    \item True Positive Rate (TPR) or Recall/Sensitivity: Measures the proportion of true positives to the number of actual positives.
    It is defined as: 

    \begin{equation}
        \text{TPR} = \frac{TP}{P} = \frac{TP}{TP + FN}
    \end{equation}

    \item Specificity: Measures the proportion of true negatives to the number of actual negatives.
    It's goal is to measure the efectiveness of the method to identify negatives.
    It is defined as:

    \begin{equation}
        \text{Specificity} = \frac{TN}{FP + TN}
    \end{equation}
\citep[p.430]{sokolova2009systematic}
\endenumerate
\todo{consider: maybe its better to plot this as a table}

\textcolor{red}{First describe the dataset, then critizise the metrics and choose them!}
\subsection{Critique of Accuracy}
Accuracy is a commonly used metric in outlier detection. \todo{adding: Papers that use accuracy as a metric}
In a comprehensive survey of 290 papers on machine learning-based anomaly detection, it was found that 112 of those papers used Accuracy \citep[p.78664]{nassif2021}. 
Nevertheless Accurarcy is not suitable for outlier detection tasks, because the datasets are imbalanced in nature. 
That means that the minority class (outlier) is much smaller than the majority class (inliers).
The impact of the outliers will be reduced when compared to the inliers.
For a small given percentage of outliers, accuracy can be estimated high by only classifying all instances as inliers.
Resulting in the given percentage of misclassified instances.
Therefore, a metric used for imbalanced datasets should consider the distributions of the classes. \citep[p.8]{branco2016survey}  


\subsection{Invariance characteristics}
To analyse the robustness of the traditional metrics, I will apply the  framework of \citet{sokolova2009systematic}.
The authers analysed different approaches to survey the invariance properties of up to 24 metrics.
Invariance is the robustness of the metrics if the confusion matrix changes.
The authers identified 7 different approaches to analyse the Invariace carakteristics of the above called metrics.
I will use three datasets wich are highly imbalanced and vary in the number of dimensions and instances.
Therefore, the following 
 
Precision, Recall and Specificity are calculated from the values of the confusion matrix and simply. 
If the threshold is changed, so must change the confusion matrix' values as well. 


Precision negates this problem by not considering the distributions but by ignoring the true negatives. 
Nevertheless, this approach has its own disadvantages. 
By ignoring the negatives the method can optimize for that metric by classifying conservatively, predicting only the most obvious outliers.


Hence, in scenarios where the cost of a false negative is high, like in the datasets I analyse, precision is not suitable as well. 


\textcolor{red}{Precision rarely operates in a vacuum. There is a fundamental trade-off governed by the classification threshold.
    -   The Problem: Increasing the threshold to improve Precision almost inevitably decreases Recall %(TP/(TP+FN))%.
    Thesis Application: In your comparison of three outlier detection methods, reporting only Precision would mask whether a method is "cherry-picking" easy outliers while failing on the complex ones. 
    This is why the F1-Score (the harmonic mean) or PR-AUC is preferred in top-tier journals.
    -   Precision does not tell you if the model is actually "learning" the structure of the data or just getting "lucky" with a specific threshold.
    The Problem: Unlike the Matthews Correlation Coefficient (MCC), which requires the model to perform well on both the minority and majority classes, Precision can be "gamed" by a model that over-predicts the majority class to minimize False Positives.
    Another disadvantage of Precision is the base rate fallacy \todo {adding: base rate fallacy and paper sokolova2009systematic}
    }


\subsection{Critique of Recall}
    - Whats the Problem?
    - Why wont I use it for my Datasets
\textcolor{red}
    -   The "Cheap" Maximization Problem
    Recall is mathematically "cheap" to maximize. An algorithm can achieve a perfect Recall of 1.0 (100\%) simply by labeling every single data point as an outlier.
    The Problem: A model that flags everything is technically perfect according to the Recall metric, but it is practically useless because it provides zero discrimination between classes.
    Recall is calculated as TP/(TP+FN). Notice that the FP (False Positive) count is entirely absent from this equation.
    -   Complete Ignorance of False Positives (Type I Errors)
    The Problem: In an econometric or business context (like fraud detection), a high-recall model might flag thousands of legitimate transactions as "fraudulent." 
    The cost of investigating these false alarms is often high, but Recall does not reflect this "noise."
    Thesis Application: If you compare your three outlier detection methods and Method A has a higher Recall than Method B, you cannot conclude Method A is better without checking if it also produces ten times as many false alarms
    -   Threshold Vulnerability
    Recall is highly dependent on the chosen classification threshold.
    The Problem: Many unsupervised outlier detection algorithms produce a continuous Anomaly Score. By lowering the threshold, you can artificially inflate Recall.
    The Result: This makes it difficult to compare algorithms fairly unless you use a threshold-independent metric like ROC-AUC or PR-AUC, which evaluate Recall across all possible thresholds.
    -   Asymmetric Information Loss
    Recall only looks at the "Positive" (Outlier) class. It tells you nothing about how well the model identifies the "Negative" (Inlier) class.
    The Problem: A model could have high Recall but be failing catastrophically at understanding the "Normal" geometric structure of your data.
    Scientific Source: Chicco  Jurman (2020) argue that because Recall ignores two quadrants of the confusion matrix (TN and FP), it provides an incomplete and potentially biased summary of model quality.
    Precision-Recall trade off

\subsection{Definition of F1-Score}
F1 score: The F1 score is the harmonic mean of
precision and recall, offering a balanced measure of model
accuracy, particularly where both false positives and false
negatives are consequential. It is defined as:
The F1 score depends on a specific threshold setting to
determine the binary classification of cases as normal or
anomaly. The choice of this threshold influences both the
precision and recall, thereby affecting the F1 score. This metric
indicates the overall robustness of anomaly detection models
in balanced class distributions [26].
    - how is it calculated
    - Where has it been used
    - Advantages and Disadvantages of F1
The F-score is evenly balanced when beta = 1. It favours precision when beta > 1,
and recall otherwise.
\textcolor{red}{
    -   The Exclusion of True Negatives (TN)
    The most significant mathematical critique of the F1 score is that it completely ignores the True Negative quadrant of the confusion matrix.
    The Problem: The formula (F1=2TP+FP+FN2TP) does not include TN. In outlier detection, the "Normal" class is the vast majority of your data.
    A metric that ignores how well you identify the normal class is technically providing an incomplete picture of the models "total" knowledge.
    -   Equal Weighting of Precision and Recall
    The F1 score is a specific case of the F$\beta$ measure where $\beta$=1, meaning it gives equal importance to Precision and Recall.
    The Problem: In many real-world scenarios (like fraud detection or medical diagnosis), the cost of a False Negative (missing a fraud) is much higher than a False Positive (a false alarm). 
    The F1 score "forces" a 1:1 trade-off that may not align with your econometric cost-benefit analysis.
    The Solution: In such cases, researchers often use F2 (weighing Recall higher) or F0.5 (weighing Precision higher).}
\subsection{Critique of Specificity}
    - Whats the Problem?
    - Why wont I use it for my Datasets
Conversely, two measures that separately estimate a classifier’s performance on different classes are sensitivity and specificity (often employed in biomedical and medical
applications, and in studies which involve image and visual data)
\textcolor{red}{
    -   The "Large Negative" Dilution Effect
    The formula for Specificity is TN/(TN+FP). In outlier detection, the number of True Negatives (TN) is usually massive (e.g., millions of normal transactions).
    The Problem: Because the denominator is dominated by the TN count, even a significant absolute number of False Positives (FP) will not noticeably move the metric.
    The Consequence: A model could produce 5,000 "false alarms" (False Positives), but if there are 1,000,000 inliers, the Specificity remains a seemingly perfect 99.5\%. 
    This hides the fact that the model is overwhelming the user with errors. \citep{saito2015precision}
    -   Complete Ignorance of the Outlier Class
    Specificity focuses entirely on the "Negative" class. It uses only the bottom row of the confusion matrix (TN and FP).
    The Problem: It tells you absolutely nothing about the model's ability to find actual outliers. 
    A model that labels everything as an inlier will achieve a perfect Specificity of 100\%, despite being a total failure at detecting anomalies.
}


\subsection{Definition of ROC and AUC}
    - how is it calculated
    - Where has it been used
    - Advantages and Disadvantages of AUC
ROC AUC (Receiver Operating Characteristic Area Un-
der Curve): ROC AUC assesses the likelihood of a random
abnormal sample being more anomalous than a normal one. It
is visualized as a curve plotting the true positive rate (TPR)
against the false positive rate (FPR), and can be mathemati-
cally expressed as an integral across all thresholds. The ROC
AUC value ranges from 0 to 1, where a higher value indicates
better model discrimination [25].
the whole two dimensional area under the entire ROC curve.
ROC curve is one of the strongest metrics used to efficiently
assess intrusion detection systems performance, and it is a
graphical tool that illustrates accuracy across FPS \citep[postfix]{nassif2021}
\todo{image of ROC}
A comprehensive evaluation of classifier performance can be obtained by the ROC 
insert ROC equation 
denotes the conditional probability that a data entry has the class label C. An
ROC curve plots the classification results from the most positive to the most negative
classification. Due to the wide use in cost/benefit decision analysis, ROC and the Area
under the Curve (AUC) apply to learning with asymmetric cost functions and imbal-
anced data sets
\textcolor{red}{critique on traditional metrics}
e argue that performance measures other than accuracy, F-score, precision, recall or
ROC do apply and can be beneficial. 
Following \citep[p.1017]{sokolova2006beyond}


\subsection{Definition of AUCPR}
    - how is it calculated
    - Where has it been used
    - Advantages and Disadvantages of AUCPR
AUCPR (Area Under the Precision-Recall Curve):
Absolutly implement into the data analysis!
AUCPR, designed for imbalanced datasets, emphasizes the mi-
nority class by plotting precision against recall across various
thresholds. It evaluates the model’s performance in scenarios
where false negatives carry more weight than false positives,
making it particularly useful in fields where the consequences
of missing true positive cases are critical [27]. Similar to
ROC AUC, AUCPR can be mathematically expressed as an
integral across the entire precision-recall curve [28]. This
integral representation provides a comprehensive assessment
of the model’s performance over various threshold levels,
reflecting its effectiveness in distinguishing between classes
in imbalanced situations. \citep{ok2024exploring}


\subsection{adding: Runtime}

